\documentclass[10pt]{article}

% ---------- Document Identity (update these only) ----------
\newcommand{\DocStage}{}
\newcommand{\DocTitle}{Your Road to Tier One SOF}
\newcommand{\ProgramName}{182-Week Civilian-to-Selection Readiness Pipeline}
\newcommand{\ProgramSubtitle}{}

% ---------- Page ----------
\usepackage[legalpaper,left=0.5in,right=0.5in,top=0.6in,bottom=0.45in]{geometry}

% ---------- Typography ----------
\usepackage[T1]{fontenc}
\usepackage[utf8]{inputenc}
\usepackage{libertinus}
\usepackage{microtype}
\usepackage{parskip}
\usepackage{ragged2e}
\usepackage{textcomp}
\AtBeginDocument{\color{black}}
\AtBeginDocument{\pagecolor{white}}
\AtBeginDocument{\justifying}

% ---------- Landscape pages ----------
\usepackage{pdflscape}

% ---------- Color ----------
\usepackage[table]{xcolor}
\definecolor{StageBlue}{HTML}{1F4E79}
\definecolor{StageBlueDark}{HTML}{163A5A}
\definecolor{LightGray}{HTML}{F2F4F7}
\definecolor{MidGray}{HTML}{D9DEE5}
\definecolor{RuleGray}{HTML}{C7CED8}
\definecolor{HeaderInk}{HTML}{000000}
\definecolor{Ink}{HTML}{000000}

% ---------- Utility ----------
\usepackage{enumitem}
\sloppy
\setlength{\emergencystretch}{2em}

% ---------- Boxes / UI ----------
\usepackage[most]{tcolorbox}

\tcbset{
  charterTitle/.style={
    enhanced,
    colback=StageBlue!4,
    colframe=StageBlue,
    boxrule=1.25pt,
    arc=3pt,
    left=16pt,right=16pt,top=14pt,bottom=14pt,
    borderline north={3.2pt}{0pt}{StageBlue},
    borderline west={4pt}{0pt}{StageBlue},
    borderline south={1.25pt}{0pt}{StageBlue},
    drop shadow={black!25!white},
  },
  charterRule/.style={
    enhanced,
    colback=LightGray,
    colframe=RuleGray,
    boxrule=0.85pt,
    arc=3pt,
    left=12pt,right=12pt,top=10pt,bottom=10pt,
    borderline west={3pt}{0pt}{StageBlue},
  },
  charterBadge/.style={
    enhanced,
    colback=StageBlue,
    coltext=white,
    colframe=StageBlueDark,
    boxrule=0.6pt,
    arc=2pt,
    left=6pt,right=6pt,top=3pt,bottom=3pt,
  }
}
\newtcbox{\Badge}{nobeforeafter,charterBadge}

% ---------- Lists ----------
\setlist[itemize]{leftmargin=1.5em, itemsep=0.25em, topsep=0.25em}
\setlist[enumerate]{leftmargin=1.8em, itemsep=0.25em, topsep=0.25em}

% ---------- TOC ----------
\usepackage{tocloft}
\setcounter{tocdepth}{3}
\setlength{\cftbeforesecskip}{3pt}
\setlength{\cftbeforesubsecskip}{2pt}
\setlength{\cftbeforesubsubsecskip}{0pt}
\renewcommand{\cftsecfont}{\bfseries}
\renewcommand{\cftsecpagefont}{\bfseries}
\renewcommand{\cftsubsecfont}{\normalfont}
\renewcommand{\cftsubsecpagefont}{\normalfont}
\renewcommand{\cftsubsubsecfont}{\normalfont}
\renewcommand{\cftsubsubsecpagefont}{\normalfont}
\setlength{\cftsecindent}{0pt}
\setlength{\cftsubsecindent}{1.25em}
\setlength{\cftsubsubsecindent}{2.8em}
\setlength{\cftsecnumwidth}{2.1em}
\setlength{\cftsubsecnumwidth}{2.8em}
\setlength{\cftsubsubsecnumwidth}{3.6em}

% Remove the built-in TOC heading so only the custom heading is shown
\makeatletter
\renewcommand{\tableofcontents}{\@starttoc{toc}}
\makeatother

% ---------- Header/Footer: page number only ----------
\usepackage{fancyhdr}
\pagestyle{fancy}
\fancyhf{}
\renewcommand{\headrulewidth}{0pt}
\renewcommand{\footrulewidth}{0pt}
\fancyfoot[C]{\thepage}

% ---------- Section formatting ----------
\usepackage{titlesec}

\titleformat{\section}
  {\Large\bfseries\color{Ink}}
  {\thesection}{0.6em}{}
  [\vspace{0.25em}\textcolor{StageBlue}{\rule{\linewidth}{0.8pt}}]

\titleformat{\subsection}
  {\large\bfseries\color{Ink}}
  {\thesubsection}{0.6em}{}

\titleformat{\subsubsection}
  {\normalsize\bfseries\color{Ink}}
  {\thesubsubsection}{0.6em}{}

\titlespacing*{\section}{0pt}{0.95em}{0.35em}
\titlespacing*{\subsection}{0pt}{0.65em}{0.25em}
\titlespacing*{\subsubsection}{0pt}{0.55em}{0.2em}

% ---------- Table engine ----------
\usepackage{tabularray}
\UseTblrLibrary{booktabs}

% ---------- tabularray: GLOBAL table treatment ----------
\SetTblrInner{
  cells = {fg=black, bg=white, valign=t},
}

% ---------- Header helpers ----------
\newcommand{\Hdr}[1]{\shortstack[c]{\strut #1\strut}}

% ---------- Lines ----------
\newcommand{\TitleLine}{\textcolor{StageBlue}{\rule{0.98\linewidth}{0.95pt}}}
\newcommand{\SubTitleLine}{\textcolor{RuleGray}{\rule{0.98\linewidth}{0.65pt}}}

% ---------- Continued on next page ----------
\newcommand{\ContinuedOnNextPageInline}{%
  \par
  \vfill
  \noindent\textcolor{RuleGray}{\rule{\linewidth}{1pt}}\vspace{-2pt}
  \noindent\hfill\textit{Continued on next page}\par\vspace{-10pt}
  \noindent\textcolor{RuleGray}{\rule{\linewidth}{1pt}}
  \clearpage
}

% ---------- PDF hyperlinks ----------
\usepackage[hidelinks]{hyperref}
\hypersetup{
  colorlinks=false,
  pdfborder={0 0 0},
  linktoc=all,
  pdftitle={\DocTitle},
  pdfauthor={\ProgramName}
}

% =========================================================
\begin{document}

\begin{landscape}

\section*{Stage 4.B — Phase 00 Timeline and Microcycle Structure}
\phantomsection
\addcontentsline{toc}{section}{Stage 4.B — Phase 00 Timeline and Microcycle Structure}

\subsection*{4.B.1 Segment Summary Table}
\phantomsection
\addcontentsline{toc}{subsection}{4.B.1 Segment Summary Table}

\begingroup
\scriptsize
\setlength{\tabcolsep}{2pt}
\renewcommand{\arraystretch}{1.12}

\begin{longtblr}{
  width=\linewidth,
  rowhead=1,
  colspec = {
    Q[l,wd=0.85in]
    Q[l,wd=0.95in]
    Q[l,wd=1.75in]
    Q[l,wd=1.05in]
    X[l,2.85]
  },
  hlines,
  vlines,
  cells = {hsep=6pt, vsep=6pt, halign=l, valign=t},
  cell{1-Z}{1-5} = {fg=black},
  row{odd} = {bg=white},
  row{even} = {bg=LightGray},
  row{1} = {bg=StageBlue!15, fg=black, font=\bfseries, halign=l, valign=t},
}
\Hdr{Segment} &
\Hdr{Weeks} &
\Hdr{Primary Focus} &
\Hdr{Work:Deload Pattern} &
\Hdr{Key Adaptation Goals} \\

Segment 1 &
Weeks 1–4 &
Habit, standardization, tolerance creation &
3:1 &
Establish 6-session execution discipline under full session-element order; stabilize footwear interface and safe surface policy; build continuous low-impact Cardio tolerance without breathless intensity; establish logging discipline with no backfill; prepare lock fields for Week 4 Deload+Test baseline battery under stable conditions. \\

Segment 2 &
Weeks 5–8 &
Aerobic minutes build, low-impact consistency &
3:1 &
Increase repeatable Cardio exposure by stability (not speed); reduce day-to-day variance in RPE and talk-test cues; maintain foot integrity under higher step throughput; preserve lock fields for chair height and incline height; execute Week 8 Mid Check battery without escalating risk posture. \\

Segment 3 &
Weeks 9–12 &
Consolidation, controlled density, repeatability &
3:1 &
Prove repeatability of low-risk patterning and Cardio under stable conditions; minimize \textbf{MODIFIED} / \textbf{INVALID} triggers; maintain stable gait under longer step-based Cardio where feasible; execute Week 12 Confirm battery under identical lock fields to Week 4 for comparability confirmation. \\

Segment 4 &
Weeks 13–16 &
Consolidation, stability proof, Phase 01 entry posture &
3:1 &
Preserve low-variance execution, stable recovery trends, and stable durability posture; enforce no-novelty posture approaching Week 16; execute Week 16 Exit Gate battery under \textbf{VALID} conditions where possible; if compromised, apply \textbf{HOLD} and reslot/extension posture without stacking. \\

\end{longtblr}

\endgroup

\subsection*{4.B.2 Segment Microcycles}
\phantomsection
\addcontentsline{toc}{subsection}{4.B.2 Segment Microcycles}

Global microcycle rules that apply to every segment:

\begin{itemize}
\item Six sessions per week are executed as the fixed cadence. Fatigue is managed by downshift and substitution, not by reducing frequency.
\item Every session follows the required element order: Safety self-check first → Warm-up → Mobility → Main work → Cardio → Cooldown → Recovery notes and required post-session notes.
\item Cardio occurs every session and is executed as a continuous block meeting the minimum standard; if continuous duration cannot be maintained safely, downgrade intensity and/or switch to a lower-risk modality to preserve continuity.
\item Step count integration is embedded inside the Cardio block: step-based Cardio is the default where feasible, and the 10,000 steps/day intent is pursued and logged as baseline activity; do not treat steps as a substitute for Cardio minutes unless explicitly logged as Cardio inside a session.
\item Cardio-primary vs non-cardio-primary posture:
\begin{itemize}
\item Cardio-primary: Cardio is the dominant stressor and remains talk-test governed.
\item Non-cardio-primary: technique and patterning are dominant; Cardio is completed as a low-cost add-on to meet the minimum standard.
\end{itemize}
\item Phase 00 remains submax and low-risk: no breathless intensity chasing, no impact running escalation, no rucking, no swim time trials, no underwater, no maximal strength.
\item No novelty posture in lead-in to Deload+Test weeks: do not change footwear system, route class/surface class, tool identity, chair height and arms policy, incline height, or scoring standard unless explicitly treated as re-baselining with validity consequences.
\end{itemize}

\end{landscape}

\begin{landscape}

\subsubsection*{Segment 1 — Weeks 1–4}
\phantomsection

\begin{samepage}
Typical Work-Week Microcycle Table for Weeks 1–3\par\vspace{0.25em}

\begingroup
\scriptsize
\setlength{\tabcolsep}{2pt}
\renewcommand{\arraystretch}{1.12}

\begin{longtblr}{
  width=\linewidth,
  rowhead=1,
  colspec = {
    Q[l,wd=0.60in]
    Q[l,wd=1.55in]
    Q[l,wd=1.55in]
    X[l,2.05]
    X[l,2.85]
  },
  hlines,
  vlines,
  cells = {hsep=6pt, vsep=6pt, halign=l, valign=t},
  cell{1-Z}{1-5} = {fg=black},
  row{odd} = {bg=white},
  row{even} = {bg=LightGray},
  row{1} = {bg=StageBlue!15, fg=black, font=\bfseries, halign=l, valign=t},
}
\Hdr{Day} &
\Hdr{Primary Focus} &
\Hdr{Secondary Focus} &
\Hdr{Example Session Theme} &
\Hdr{Notes — Intensity and Duration Caps} \\

Day 1 &
Execution discipline and stable gait &
\textbf{DURABILITY} baseline capture &
Cardio-primary: step-based walking where feasible; stable surface class; Cardio completed as continuous block; technique-first main work patterning &
Cardio: continuous, minimum standard met; RPE 2–4; talk test full sentences. Steps posture pursued through Cardio intent and daily activity; log steps count. No terrain novelty; stop and substitute if gait changes or hotspots appear. \\

Day 2 &
Mobility continuity &
Core bracing baseline &
Non-cardio-primary: mobility emphasis; bracing/trunk control; Cardio completed as low-cost add-on &
Cardio: continuous, minimum standard met; RPE 2–3; talk test full sentences. Technique gate: stop before form degradation. Any pain-driven compensation triggers immediate regression and \textbf{MODIFIED} tagging. \\

Day 3 &
Aerobic tolerance initiation &
Foot interface governance &
Cardio-primary: step-based walking emphasis; conservative environment selection; Cardio completed as continuous block &
Cardio: continuous, minimum standard met; RPE 3–4; talk test full sentences to short sentences. Foot inspection before/after; any hotspot trend triggers next-session downshift. Indoor substitute permitted if surface is unsafe. \\

Day 4 &
Patterning rehearsal for lock fields &
Mobility reinforcement &
Non-cardio-primary: chair and incline setup rehearsal as patterning only; Cardio completed as low-cost add-on &
Cardio: continuous, minimum standard met; RPE 2–4; talk test full sentences. Record candidate lock fields (chair height, arms policy, incline height) in notes without treating the day as a test. \\

Day 5 &
Recovery-throughput execution &
\textbf{RECOVERY} marker stability &
Cardio-primary: lowest-risk modality; Cardio completed as continuous block; cooldown and recovery notes emphasized &
Cardio: continuous, minimum standard met; RPE 2–3; talk test full sentences. Preserve low variance; do not “push harder” because it feels easy. Log sleep/soreness/fatigue/pain and steps. \\

Day 6 &
Consolidation and pre-window stabilization &
Documentation quality check &
Non-cardio-primary: technique-first patterning; Cardio completed as low-cost add-on &
Cardio: continuous, minimum standard met; RPE 2–4; talk test full sentences. Treat as stabilization for Week 4: no footwear/route/tool changes; confirm lock-field feasibility. \\

\end{longtblr}

\endgroup
\end{samepage}

Day 7 — Non-Session Recovery Movement Logging\par\vspace{0.25em}

\begin{itemize}
\item Low-intensity recovery intent: maintain movement without creating new fatigue cost or soreness spikes.
\item Mobility and light movement categories: gentle mobility, easy walking where safe, and downregulation breathing as tolerated.
\item Required logging and review actions:
\begin{itemize}
\item Record sleep hours, sleep quality, soreness 0–10, fatigue 0–10, pain 0–10 with location when present, and foot hotspot status.
\item Review the prior six sessions for completeness of required elements and documentation.
\item Confirm Cardio was recorded as a continuous block each session and that steps were logged daily.
\end{itemize}
\item Step-based activity guidance:
\begin{itemize}
\item Steps posture is maintained as baseline activity; if environment or foot integrity is compromised, use the lowest-risk movement available and document constraints.
\end{itemize}
\end{itemize}

Deload+Test Week Modifiers — Week 4\par\vspace{0.25em}

Deload+Test week posture at the scheduling level:

\begin{itemize}
\item Reduced volume density and reduced intensity exposure relative to work weeks, while preserving the six-session rhythm, preserving complete session-element order, and preserving Cardio as a continuous block meeting the minimum standard.
\item Testing is staged to protect validity and reduce risk; no day-by-day test schedule is created in this deliverable.
\item If any scheduled item cannot meet \textbf{VALID} conditions due to illness, pain drift, missing lock fields, environmental uncertainty, or tool change, treat it as \textbf{INVALID} for gate use and reslot to the next Deload+Test window.
\end{itemize}

Week 4 Deload+Test Baseline Battery:

\begin{itemize}
\item \textbf{SMB}
\item \textbf{MOB} Bundle: \textbf{MOB-MOV-SCREEN}; \textbf{MOB-ANKDF-K2W}; \textbf{MOB-SHFLEX-WALL}; \textbf{MOB-STS-QUAL}
\item \textbf{ME-L-STS-30}
\item \textbf{ME-UP-PU-INC}
\item \textbf{CAL-PLANK}
\item \textbf{AER-WALK-30}
\end{itemize}

Comparability locks required at governance level:

\begin{itemize}
\item Footwear and sock system locked.
\item Route class and surface class or indoor substitute policy locked for walk-based monitoring.
\item Chair height and arms policy locked for the sit-to-stand instrument.
\item Incline push-up height locked for incline push instrument.
\item Walk monitor route conditions locked: surface class and environment notes recorded consistently.
\item Tool standard consistency: timer method and measurement method are stable; any change is documented and treated as a comparability break.
\item Mode logging posture: any required mode fields are recorded explicitly when applicable.
\end{itemize}

Validity and logging reminder:

\begin{itemize}
\item Validity and logging standards are governed by Stage 1.F and Stage 2.E. Missing lock fields or missing required documentation make evidence inadmissible for gate decisions.
\end{itemize}

\subsubsection*{Segment 2 — Weeks 5–8}
\phantomsection

\begin{samepage}
Typical Work-Week Microcycle Table for Weeks 5–7\par\vspace{0.25em}

\begingroup
\scriptsize
\setlength{\tabcolsep}{2pt}
\renewcommand{\arraystretch}{1.12}

\begin{longtblr}{
  width=\linewidth,
  rowhead=1,
  colspec = {
    Q[l,wd=0.60in]
    Q[l,wd=1.55in]
    Q[l,wd=1.55in]
    X[l,2.05]
    X[l,2.85]
  },
  hlines,
  vlines,
  cells = {hsep=6pt, vsep=6pt, halign=l, valign=t},
  cell{1-Z}{1-5} = {fg=black},
  row{odd} = {bg=white},
  row{even} = {bg=LightGray},
  row{1} = {bg=StageBlue!15, fg=black, font=\bfseries, halign=l, valign=t},
}
\Hdr{Day} &
\Hdr{Primary Focus} &
\Hdr{Secondary Focus} &
\Hdr{Example Session Theme} &
\Hdr{Notes — Intensity and Duration Caps} \\

Day 1 &
Cardio consistency build &
\textbf{DURABILITY} stability monitoring &
Cardio-primary: step-based walking where feasible; Cardio continuous block; stable surface class &
Cardio: continuous, minimum standard met; RPE 3–4; talk test full sentences to short sentences. Avoid adding speed targets; increase only by stability. Document foot and gait stability. \\

Day 2 &
Mobility and joint-control continuity &
Patterning stability &
Non-cardio-primary: mobility emphasis; chair and incline patterning at low strain; Cardio add-on &
Cardio: continuous, minimum standard met; RPE 2–3; talk test full sentences. No novelty; lock fields remain stable for Week 8. \\

Day 3 &
Aerobic repeatability &
\textbf{RECOVERY} stability &
Cardio-primary: controlled, repeatable Cardio exposure; minimize intensity variance &
Cardio: continuous, minimum standard met; RPE 3–4; talk test stable. If sleep streak trigger appears, downshift intensity and/or substitute modality without reducing frequency. \\

Day 4 &
Trunk endurance continuity &
Mobility reinforcement &
Non-cardio-primary: trunk control emphasis; Cardio add-on &
Cardio: continuous, minimum standard met; RPE 2–4; talk test full sentences. Terminate patterning on mechanics change or pain >5/10. \\

Day 5 &
Recovery-throughput day &
Foot interface check &
Cardio-primary: lowest-risk steady modality; Cardio continuous block &
Cardio: continuous, minimum standard met; RPE 2–3; talk test full sentences. This day exists to prevent fatigue stacking and protect Week 8 validity. \\

Day 6 &
Pre-window stabilization &
Tool and environment stability &
Non-cardio-primary: minimal strain; Cardio add-on &
Cardio: continuous, minimum standard met; RPE 2–4; talk test full sentences. Confirm lock fields remain unchanged; if they must change, treat Week 8 comparability as compromised and document. \\

\end{longtblr}

\endgroup
\end{samepage}

Day 7 — Non-Session Recovery Movement Logging\par\vspace{0.25em}

\begin{itemize}
\item Low-intensity recovery intent: maintain movement and mobility without adding fatigue cost.
\item Mobility and light movement categories: gentle mobility, easy walking where safe, and downregulation.
\item Required logging and review actions:
\begin{itemize}
\item Confirm steps are logged daily and support the 10,000 steps/day intent.
\item Review whether any sessions required \textbf{MODIFIED} or \textbf{INVALID} tagging and document root cause (pain, environment, time, equipment, logging).
\item Confirm lock fields remain stable for Week 8 Mid Check (footwear interface, route policy, chair height and arms policy, incline height, timer method).
\end{itemize}
\item Step-based activity guidance:
\begin{itemize}
\item Maintain baseline steps where feasible; if environment hazards exist, select the safest indoor movement available and document constraints.
\end{itemize}
\end{itemize}

Deload+Test Week Modifiers — Week 8\par\vspace{0.25em}

Deload+Test week posture at the scheduling level:

\begin{itemize}
\item Reduced volume density and reduced intensity exposure relative to work weeks, while preserving six sessions, preserving session-element order, and preserving Cardio as a continuous block meeting the minimum standard.
\item Testing is staged to protect validity and avoid stacking high-cost exposures; no day-by-day schedule is created here.
\item If validity cannot be preserved, treat impacted items as \textbf{INVALID} for gate use and reslot.
\end{itemize}

Week 8 Deload+Test Mid Check Battery:

\begin{itemize}
\item \textbf{SMB}
\item \textbf{MOB} Bundle
\item \textbf{ME-L-STS-30}
\item \textbf{ME-UP-PU-INC}
\item \textbf{CAL-PLANK}
\item \textbf{AER-WALK-015} mode logged
\end{itemize}

Comparability locks required at governance level:

\begin{itemize}
\item Preserve Week 4 lock fields:
\begin{itemize}
\item footwear and sock system locked,
\item route class and surface class or indoor substitute policy locked,
\item chair height and arms policy locked,
\item incline height locked,
\item tool standard consistency preserved.
\end{itemize}
\item \textbf{AER-WALK-015} mode logging is mandatory and remains unresolved pending Stage 8 review because no reconciled Stage 2 \textbf{ID} is assigned here.
\end{itemize}

Validity and logging reminder:

\begin{itemize}
\item Phase 00 does not accept “close enough” conditions for gate-grade evidence. Any missing lock field or missing required documentation makes the item inadmissible for gate decisions.
\end{itemize}

\newpage

\subsubsection*{Segment 3 — Weeks 9–12}
\phantomsection

\begin{samepage}
Typical Work-Week Microcycle Table for Weeks 9–11\par\vspace{0.25em}

\begingroup
\scriptsize
\setlength{\tabcolsep}{2pt}
\renewcommand{\arraystretch}{1.12}

\begin{longtblr}{
  width=\linewidth,
  rowhead=1,
  colspec = {
    Q[l,wd=0.60in]
    Q[l,wd=1.55in]
    Q[l,wd=1.55in]
    X[l,2.05]
    X[l,2.85]
  },
  hlines,
  vlines,
  cells = {hsep=6pt, vsep=6pt, halign=l, valign=t},
  cell{1-Z}{1-5} = {fg=black},
  row{odd} = {bg=white},
  row{even} = {bg=LightGray},
  row{1} = {bg=StageBlue!15, fg=black, font=\bfseries, halign=l, valign=t},
}
\Hdr{Day} &
\Hdr{Primary Focus} &
\Hdr{Secondary Focus} &
\Hdr{Example Session Theme} &
\Hdr{Notes — Intensity and Duration Caps} \\

Day 1 &
Aerobic repeatability under low variance &
\textbf{DURABILITY} drift prevention &
Cardio-primary: step-based walking where feasible; Cardio continuous block; stable surface class &
Cardio: continuous, minimum standard met; RPE 3–4; talk test stable. No novelty in footwear or route class; reslot if hazards appear. \\

Day 2 &
Mobility continuity &
Patterning quality &
Non-cardio-primary: mobility emphasis; conservative patterning; Cardio add-on &
Cardio: continuous, minimum standard met; RPE 2–3; talk test full sentences. Preserve joint-history constraints; no end-range forcing. \\

Day 3 &
Tolerance consolidation &
Foot integrity governance &
Cardio-primary: steady Cardio; gait monitoring; Cardio continuous block &
Cardio: continuous, minimum standard met; RPE 3–4; talk test stable. Any hotspot trend triggers immediate downshift next session. \\

Day 4 &
Core and posture control &
Mobility reinforcement &
Non-cardio-primary: trunk control emphasis; Cardio add-on &
Cardio: continuous, minimum standard met; RPE 2–4; talk test full sentences. Terminate any exposure that provokes mechanics change or unusual breathlessness. \\

Day 5 &
Recovery-throughput day &
\textbf{RECOVERY} marker stability &
Cardio-primary: lowest-risk steady modality; Cardio continuous block &
Cardio: continuous, minimum standard met; RPE 2–3; talk test full sentences. Preserve sleep and soreness stability; do not introduce intensity spikes. \\

Day 6 &
Pre-window stabilization for confirm week &
Lock-field verification &
Non-cardio-primary: minimal strain; Cardio add-on &
Cardio: continuous, minimum standard met; RPE 2–4; talk test full sentences. Confirm identical lock fields to Week 4 are feasible for Week 12 Confirm. \\

\end{longtblr}

\endgroup
\end{samepage}

Day 7 — Non-Session Recovery Movement Logging\par\vspace{0.25em}

\begin{itemize}
\item Low-intensity recovery intent: preserve readiness and prevent fatigue carryover into Week 12.
\item Mobility and light movement categories: gentle mobility, easy walking where safe, and downregulation.
\item Required logging and review actions:
\begin{itemize}
\item Verify lock fields remain identical to Week 4 for Week 12 Confirm (footwear, route policy, chair and incline heights, timer method).
\item Review trend stability in sleep, soreness, fatigue, pain location, and foot hotspots; plan conservative downshift if drift is present.
\item Confirm steps are logged daily and reflect the 10,000 steps/day intent as baseline activity.
\end{itemize}
\end{itemize}

Deload+Test Week Modifiers — Week 12\par\vspace{0.25em}

Deload+Test week posture at the scheduling level:

\begin{itemize}
\item Reduced density and reduced intensity compared to work weeks while preserving six sessions, preserving session-element order, and preserving Cardio as a continuous block meeting the minimum standard.
\item Testing is staged to protect validity; no day-by-day schedule is created here.
\item Confirm week exists to prove comparability under identical conditions; it does not authorize escalation.
\end{itemize}

Week 12 Deload+Test Confirm Battery:

\begin{itemize}
\item Repeat the Week 4 battery under identical lock fields, including the same locked aerobic monitor selection:
\begin{itemize}
\item \textbf{SMB}
\item \textbf{MOB} Bundle: \textbf{MOB-MOV-SCREEN}; \textbf{MOB-ANKDF-K2W}; \textbf{MOB-SHFLEX-WALL}; \textbf{MOB-STS-QUAL}
\item \textbf{ME-L-STS-30}
\item \textbf{ME-UP-PU-INC}
\item \textbf{CAL-PLANK}
\item \textbf{AER-WALK-30} (Requires Stage 8 review - no reconciled Stage 2 \textbf{ID} assigned)
\end{itemize}
\end{itemize}

Comparability locks required at governance level:

\begin{itemize}
\item All Week 4 lock fields must remain identical:
\begin{itemize}
\item footwear and sock system locked,
\item route class and surface class or indoor substitute policy locked,
\item chair height and arms policy locked,
\item incline push-up height locked,
\item walk monitor route conditions locked,
\item tool standard consistency maintained.
\end{itemize}
\item If any lock field must change:
\begin{itemize}
\item document the change,
\item treat comparability as compromised for gate interpretation,
\item reslot any gate-grade interpretation to the next valid window rather than forcing conclusions.
\end{itemize}
\end{itemize}

Validity and logging reminder:

\begin{itemize}
\item Missing lock fields or missing required documentation make evidence inadmissible for gate use. Confirm week is not a permission to accept compromised evidence.
\end{itemize}

\newpage

\subsubsection*{Segment 4 — Weeks 13–16}
\phantomsection

\begin{samepage}
Typical Work-Week Microcycle Table for Weeks 13–15\par\vspace{0.25em}

\begingroup
\scriptsize
\setlength{\tabcolsep}{2pt}
\renewcommand{\arraystretch}{1.12}

\begin{longtblr}{
  width=\linewidth,
  rowhead=1,
  colspec = {
    Q[l,wd=0.60in]
    Q[l,wd=1.55in]
    Q[l,wd=1.55in]
    X[l,2.05]
    X[l,2.85]
  },
  hlines,
  vlines,
  cells = {hsep=6pt, vsep=6pt, halign=l, valign=t},
  cell{1-Z}{1-5} = {fg=black},
  row{odd} = {bg=white},
  row{even} = {bg=LightGray},
  row{1} = {bg=StageBlue!15, fg=black, font=\bfseries, halign=l, valign=t},
}
\Hdr{Day} &
\Hdr{Primary Focus} &
\Hdr{Secondary Focus} &
\Hdr{Example Session Theme} &
\Hdr{Notes — Intensity and Duration Caps} \\

Day 1 &
Stability under routine &
\textbf{DURABILITY} drift prevention &
Cardio-primary: step-based walking where feasible; Cardio continuous block; stable surface class &
Cardio: continuous, minimum standard met; RPE 3–4; talk test stable. No novelty; treat this as execution proof, not progression chasing. \\

Day 2 &
Mobility continuity &
Recovery stabilization &
Non-cardio-primary: mobility emphasis; minimal strain patterning; Cardio add-on &
Cardio: continuous, minimum standard met; RPE 2–3; talk test full sentences. Protect joint-history constraints; no end-range forcing. \\

Day 3 &
Aerobic continuity with gait stability &
Foot interface governance &
Cardio-primary: steady Cardio; gait monitoring; Cardio continuous block &
Cardio: continuous, minimum standard met; RPE 3–4; talk test stable. Downshift immediately if gait changes or hotspots worsen. \\

Day 4 &
Lock-field rehearsal without testing &
Core and posture continuity &
Non-cardio-primary: rehearsal of chair/incline setups as patterning only; Cardio add-on &
Cardio: continuous, minimum standard met; RPE 2–4; talk test full sentences. Confirm Week 16 lock fields are feasible without changes. \\

Day 5 &
Recovery-throughput day &
Documentation quality &
Cardio-primary: lowest-risk steady modality; Cardio continuous block &
Cardio: continuous, minimum standard met; RPE 2–3; talk test full sentences. Preserve low variance; emphasize recovery notes completeness. \\

Day 6 &
Pre-gate stabilization &
No-novelty enforcement &
Non-cardio-primary: minimal strain; Cardio add-on &
Cardio: continuous, minimum standard met; RPE 2–3; talk test full sentences. Enforce no novelty for Week 16; document any constraint that would compromise validity and plan reslot posture. \\

\end{longtblr}

\endgroup
\end{samepage}

Day 7 — Non-Session Recovery Movement Logging\par\vspace{0.25em}

\begin{itemize}
\item Low-intensity recovery intent: preserve readiness for Week 16 and avoid new soreness spikes.
\item Mobility and light movement categories: gentle mobility, easy walking where safe, and downregulation.
\item Required logging and review actions:
\begin{itemize}
\item Confirm the prior work-week maintained complete session element order and complete documentation.
\item Confirm sleep/soreness/fatigue/pain/foot logs show stability and no emerging risk flags.
\item Confirm steps are logged daily and reflect the 10,000 steps/day intent as baseline activity.
\item Confirm Week 16 lock fields are ready and stable (footwear, route policy, chair/incline heights, tool identity).
\end{itemize}
\end{itemize}

Deload+Test Week Modifiers — Week 16\par\vspace{0.25em}

Deload+Test week posture at the scheduling level:

\begin{itemize}
\item Reduced density and reduced intensity compared to work weeks while preserving six sessions, preserving session-element order, and preserving Cardio as a continuous block meeting the minimum standard.
\item Testing is staged to protect validity and reduce risk; no day-by-day test schedule is created here.
\item If any scheduled test cannot meet \textbf{VALID} conditions due to illness, pain drift, missing lock fields, environmental uncertainty, or tool change, treat it as \textbf{INVALID} for gate use and reslot per Stage 3.
\end{itemize}

Week 16 Deload+Test Exit Gate Battery:

\begin{itemize}
\item Required:
\begin{itemize}
\item \textbf{SMB} (Requires Stage 8 review - no reconciled Stage 2 \textbf{ID} assigned)
\item \textbf{MOB} Bundle
\item \textbf{ME-L-STS-30}
\item \textbf{ME-UP-PU-INC}
\item \textbf{CAL-PLANK}
\item \textbf{AER-WALK-015} (Requires Stage 8 review - no reconciled Stage 2 \textbf{ID} assigned)
\end{itemize}
\item Optional baseline-only items if strictly viable and risk-neutral:
\begin{itemize}
\item \textbf{C4} — Push-Ups — 2-Minute Timed, Strict Standard (optional)
\item \textbf{C5} — Sit-Ups — 2-Minute Timed, Standard (optional)
\end{itemize}
\end{itemize}

Comparability locks required at governance level:

\begin{itemize}
\item Footwear and sock system locked; any change is documented and treated as a comparability break.
\item Route class and surface class or indoor substitute policy locked; unsafe conditions trigger indoor substitution only when comparability can be preserved and logged; otherwise reslot.
\item Chair height and arms policy locked for the sit-to-stand instrument.
\item Incline push-up height locked.
\item Tool standard consistency: timer method and measurement method are stable; late or reconstructed entries are inadmissible for gate decisions unless objectively verifiable and marked as verified late entry.
\item Mode logging for \textbf{AER-WALK-015} is mandatory and must match prior mode posture for comparability.
\end{itemize}

Validity and logging reminder:

\begin{itemize}
\item Validity tagging and admissibility are governed by Stage 1.F and Stage 2.E. Missing lock fields or missing required documentation makes evidence inadmissible for gate decisions.
\end{itemize}

\end{landscape}

\end{document}